\documentclass[12pt]{article}
\newcommand{\level}{Grades 2--3}
\usepackage[letterpaper,top=0.45in,bottom=0.52in,left=0.5in,right=0.5in,
  includehead,includefoot,headheight=18pt,headsep=0.22in,footskip=0.34in]{geometry}
\usepackage[T1]{fontenc}
\usepackage[default]{sourcesanspro}
\usepackage{tikz}
\usetikzlibrary{arrows.meta}
\usepackage{fancyhdr}
\usepackage{array}
\pagestyle{fancy}
\fancyhf{}
\fancyhead[L]{Week 1 / Tiling / \level}
\fancyfoot[L]{Bellingham Math Circle / Week 1}
\fancyfoot[R]{\thepage}
\renewcommand{\headrulewidth}{0pt}
\renewcommand{\footrulewidth}{0pt}
\setlength{\parindent}{0pt}
\setlength{\parskip}{0pt}
\raggedright
\newcommand{\prob}[1]{\textbf{Problem #1:}}
\newcommand{\fig}[1]{\input{figs/#1.tex}}
\newcommand{\rarrow}{\tikz[baseline=-3pt]\draw[-{Stealth[length=9pt,width=8pt]},line width=1.3pt] (0,0)--(0.42in,0);}
\newcommand{\lrarrow}{\tikz[baseline=-3pt]\draw[{Stealth[length=9pt,width=8pt]}-{Stealth[length=9pt,width=8pt]},line width=1.3pt] (0,0)--(0.6in,0);}
\newcommand{\blank}[1][0.9in]{\rule{#1}{0.6pt}}
\newcommand{\green}{\raisebox{-1pt}{\fig{greenicon}}}
\newsavebox{\figbox}
\newcommand{\labfig}[2]{\sbox{\figbox}{\fig{#2}}\raisebox{\dimexpr\ht\figbox-1.6ex\relax}{\makebox[0.3in][l]{\textbf{#1}}}\usebox{\figbox}}

\newcommand{\lab}[1]{\makebox[0.35in][l]{\textbf{#1}}}
\newcommand{\answerlines}[1]{\begin{minipage}[b]{2.7in}\foreach \i in {1,...,#1}{\rule{0pt}{0.42in}\rule{2.7in}{0.5pt}\par}\end{minipage}}
\begin{document}

In every problem, blocks go on a board along its lines, and they may not overlap or stick out past its edge.

\bigskip
\prob{1} Which of these boards can be covered with blue blocks only? Cover each board that you can, and draw lines on it to show where your blocks went. Put an X on each board that cannot be covered.

\vspace{0.4in}
\centering
\labfig{A}{rhomb2}\hspace{0.8in}\labfig{B}{trap31}

\vspace{0.55in}
\labfig{C}{h211}\hspace{0.8in}\labfig{D}{bumptri}

\vspace{0.55in}
\labfig{E}{trap2}

\raggedright
\newpage

\centering
\labfig{F}{para32}

\vspace{0.7in}
\labfig{G}{bowtie}

\raggedright
\newpage

\prob{2} Cover each triangle with blue blocks and green triangles, using as few green triangles as you can. Write how many green triangles you used.

\vspace{0.35in}
\centering
\begin{tabular}{c@{\hspace{1.0in}}c}
\labfig{A}{tri2} & \labfig{B}{tri3}\\[12pt]
green triangles: \blank & green triangles: \blank
\end{tabular}

\vspace{0.5in}
\labfig{C}{tri4}

\vspace{10pt}
green triangles: \blank

\raggedright
\newpage

\centering
\labfig{D}{tri5}

\vspace{10pt}
green triangles: \blank

\raggedright
\newpage

\prob{3} Each board below is a big hexagon with two gray holes, and no block may cover a hole. Which of these boards can be covered with blue blocks? For each board that cannot be covered, explain why.

\vspace{0.35in}
\labfig{A}{holes_0}\hfill\answerlines{6}

\vspace{0.45in}
\labfig{B}{holes_1}\hfill\answerlines{6}

\newpage

\labfig{C}{holes_2}\hfill\answerlines{6}

\vspace{0.6in}
\labfig{D}{holes_3}\hfill\answerlines{6}

\newpage

\prob{4} Could triangle C in Problem 2 be covered with blue blocks and fewer green triangles than you used? Explain.

\vspace{3.6in}
\prob{5} Imagine a giant triangle board like the ones in Problem 2, with sides 10 inches long. If it is covered with blue blocks and green triangles, what is the fewest number of green triangles that can be used? Explain how you know. Then answer the same question for a triangle with sides 100 inches long.

\newpage

\prob{6} Use purple blocks and green triangles, with no blue blocks. Cover each triangle in Problem 2 again, using as few green triangles as you can, and fill in the table. Is there a triangle where purple blocks need fewer green triangles than blue blocks? Explain.

\vspace{0.4in}
\centering
\renewcommand{\arraystretch}{2.2}
\begin{tabular}{|l|>{\centering\arraybackslash}p{0.75in}|>{\centering\arraybackslash}p{0.75in}|>{\centering\arraybackslash}p{0.75in}|>{\centering\arraybackslash}p{0.75in}|}
\hline
Triangle & A & B & C & D\\ \hline
Green triangles with blue blocks & & & & \\ \hline
Green triangles with purple blocks & & & & \\ \hline
\end{tabular}

\raggedright
\newpage

\prob{7} On the grid below, draw a board made of exactly 12 small triangles joined along their edges, so that every way of covering it with blue blocks and green triangles uses at least 4 green triangles.

\vspace{0.4in}
\centering
\fig{gridpaper}

\end{document}
